\documentclass[11pt]{article}
\usepackage[a4paper,margin=27mm]{geometry}
\usepackage{amsmath,amssymb,mathtools}
\usepackage{fontspec,unicode-math}
\setmainfont{Latin Modern Roman}
\setmathfont{Latin Modern Math}
\usepackage{xurl,hyperref}
\hypersetup{colorlinks=true,urlcolor=blue,linkcolor=blue}
\newcommand{\sref}[2]{\hyperlink{source:#1:#2}{[S#2]}}
\newcommand{\cataloguescope}[1]{\paragraph{Mathematical question.} #1}
\begin{document}
\section{Weak Infinitary Comprehension versus Weak Choice}\label{statement}

\subsection{Definitions and mathematical statement}

\paragraph{Shared notation.}$\mathbb N=\{1,2,\ldots\}$, and $\mathbb Z,\mathbb Q,\mathbb R,\mathbb C$ denote the integers, rationals, reals and complex numbers. Any other conventions are specified locally.


Work over $\mathsf{RCA}_0$, the subsystem of second-order arithmetic with $\Delta^0_1$ comprehension and $\Sigma^0_1$ induction; all structures and orders are countable and coded by sets of integers. $\mathsf{ACA}_0$ adds comprehension for every arithmetical formula (set parameters allowed), and $\mathsf{ATR}_0$ adds arithmetical transfinite recursion along every coded well-order. Reverse-mathematical equivalence means provability of both implications over this base. An $L_{\omega_1,\omega}$ sentence of arithmetic is coded by a well-founded syntax tree allowing countable conjunctions and disjunctions, alongside ordinary finite first-order operations. A valuation assigns truth values to all subformulas (and assignments to free number variables) satisfying their logical rules. Weak infinitary comprehension asserts: if $(\varphi_n)_{n\ge0}$ is a sequence of these sentences, each individually admitting a valuation, there is a set $X$ with $n\in X$ exactly when $\varphi_n$ evaluates true. Weak $\Sigma^1_1$ choice is the scheme \[\forall n\,\exists!X\,\psi(n,X)\ \Longrightarrow\ \exists Y\,\forall n\,\psi(n,Y^{[n]}),\] for every arithmetical $\psi$ with parameters, where $Y^{[n]}=\{k:\langle n,k\rangle\in Y\}$ using a fixed integer pairing function.
\paragraph{Statement recorded in the supplied source.}
Are weak $L_{\omega_1,\omega}$ comprehension and weak $\Sigma^1_1$ choice equivalent over $\mathsf{RCA}_0$? \sref{statement}{3} \sref{statement}{2}
\subsection{Short English statement}
Does collecting truth values of an individually evaluable sequence of infinitary arithmetic sentences require exactly weak unique-set choice?
\subsection{Sources}
\begin{description}
\item[\hypertarget{source:statement:1}{S1}] ULAM.AI and the UnsolvedMath Contributors, UnsolvedMath: A Curated Collection of Open Mathematics Problems, catalogue entry. CC BY 4.0. \url{https://huggingface.co/datasets/ulamai/UnsolvedMath}.
\item[\hypertarget{source:statement:2}{S2}] Antonio Montalbán, Open Questions in Reverse Mathematics, author preprint (17 August 2010), Question 31, Section 5, p. 12. \url{https://math.berkeley.edu/~antonio/papers/questionsRM.pdf}.
\item[\hypertarget{source:statement:3}{S3}] Antonio Montalbán, Indecomposable Linear Orderings and Hyperarithmetic Analysis (2006), Section 1.6, definition of weak $\Sigma^1_1$ choice. \url{https://math.berkeley.edu/~antonio/papers/indec.pdf}.
\end{description}
\label{end:statement}
\end{document}
